\section{Suspension Geometry}

\begin{multicols}{2}
The suspension geometry serves as the critical dynamic interface between the 800 kg total vehicle mass, the aerodynamic aerodynamic platform, and the physical tire contact patches. For the 2026 regulations, the suspension architecture must forcefully extract mechanical grip from the 350 kW electrical torque delivery while rigidly preserving the 32 mm front and 30 mm rear ride heights mandated by the underfloor aerodynamics. This section defines the macroscopic kinematic load paths and dynamic frequency responses governing the vehicle.

\subsection{Wishbone Geometry}

Both the front and rear axles utilize a non-parallel, unequal-length double wishbone structural configuration. The upper wishbones are geometrically shorter than the lower wishbones. This fundamental inequality generates a progressive negative camber gain during suspension bump (compression), actively tilting the top of the tire inboard to counteract external chassis roll and aggressively maintain the maximum area of the tire contact patch under severe lateral loads.

The inboard suspension pickup points are structurally anchored to the carbon fiber monocoque at the front axle and dynamically mounted to the stressed 590 mm gearbox casing at the rear. The outboard pickups articulate via spherical bearings onto the titanium (Ti-6Al-4V) uprights. The wishbones themselves are constructed from aerodynamically-profiled Carbon Fiber Reinforced Polymer (CFRP), chosen for its extraordinarily high specific stiffness to practically eliminate bending deflection while minimizing the critical unsprung corner mass. 

The geometry is optimized longitudinally for the mandated 1600 mm front track and 1620 mm rear track. The kinematic convergence point of the upper and lower wishbone extensions in the transverse plane defines the Instant Center (IC). The location of the IC directly dictates both the rate of camber gain and the absolute height of the roll center—metrics that permanently define the vehicle's transient roll stiffness distribution.

\subsection{Pullrod Actuation Strategy}

Departing from historical conventions, both the front and rear axles mandate pullrod actuation architectures. At the front, the pullrod physically connects the upper outboard wishbone juncture to a lower, inboard mechanical rocker. This geometry fundamentally lowers the physical placement of the heavy torsion springs, rotary dampers, and inerter assemblies deep into the chassis floor, aggressively supporting the stringent 270 mm overall Center of Gravity (CoG) aerodynamic target.

At the rear, the pullrod configuration remains consistent with modern design trends, proving highly compatible with the dense rear-mid aerodynamic packaging constraints surrounding the V6 Internal Combustion Engine (ICE) and longitudinal gearbox. By pulling upward on an inboard rocker mounted low on the transmission casing, the geometric load path accurately transfers the immense vertical forces from the tire contact patch directly into the centralized, low-slung spring/damper assembly.

\subsection{Camber Settings}

To preemptively counteract severe centrifugal tire carcass deformation under multi-G cornering speeds, static negative camber is applied to all four corners. 

Static camber targets are set aggressively: the front axle ranges from -2.5$^\circ$ to -3.5$^\circ$, while the traction-limited rear axle utilizes a slightly shallower -1.5$^\circ$ to -2.0$^\circ$ to maintain a flatter longitudinal contact patch for the 350 kW Motor Generator Unit-Kinetic (MGU-K) deployment. The dynamic camber gain rate ($\frac{\partial\gamma}{\partial z}$) is governed by the instant center geometry derived previously. For traditional circuit configurations, the static camber setup is perfectly symmetric across the left and right planes, contrasting heavily with the highly asymmetric setups utilized in strictly oval-based motorsport formats.

\subsection{Toe Settings}

Ackermann steering geometry is augmented by baseline static toe settings. The front axle utilizes a deliberate static toe-out configuration of $\sim$-0.05$^\circ$. This divergent angular setup intentionally induces a minor slip angle in the inner tire during initial turn-in, drastically improving the yaw response and steering agility of the vehicle on corner entry. 

Conversely, the rear axle employs a slight static toe-in geometry of $\sim$+0.1$^\circ$. The convergent rear tires fundamentally push against each other symmetrically, vastly increasing high-speed straight-line stability—a critical priority for maintaining aerodynamic poise at $v > 300$ km/h. However, engineers must fiercely account for compliance steer: the dynamic, unintended change in toe angles generated by lateral and longitudinal forces elastically defeating the suspension bush stiffness. The compliance toe change ($\Delta_{toe}$) is mathematically derived by multiplying the lateral force ($F_y$) by the localized compliance matrix [$C$].

\subsection{Caster and Kingpin Inclination}

The steering axis geometry is governed by combined Caster and Kingpin Inclination (KPI) angles. A pronounced caster angle (10-12$^\circ$) inclines the steering axis rearward. The resultant mechanical trail ($t_m$) is derived as:

\begin{equation} \label{eq:mechanical_trail}
t_m = r_{wheel} \cdot \tan(\theta_{caster})
\end{equation}

where $r_{wheel}$ is the dynamic loaded tire radius and $\theta_{caster}$ is the caster angle. A high caster angle aggressively forces the outboard tire into deeper negative camber proportionally to the steering wheel input, maximizing mid-corner grip, albeit at the severe expense of driver physical steering effort.

The Kingpin Inclination (8-10$^\circ$) inclines the upper steering pivot inboard. The lateral ground-plane distance between the steering axis intersection and the true center of the tire contact patch defines the scrub radius ($r_{scrub}$):

\begin{equation} \label{eq:scrub_radius}
r_{scrub} = r_{wheel} \cdot \tan(\theta_{kpi})
\end{equation}

where $\theta_{kpi}$ is the KPI angle. Minimizing the scrub radius is absolutely paramount to reducing violent asymmetrical steering wheel kickback transmitted to the driver under heavy, single-wheel braking lockups. The combined geometric synthesis of caster and KPI defines the total pneumatic trail, the primary source of tactile steering feedback communicating terminal grip limits to the driver.

\subsection{Anti-Dive and Anti-Squat Geometry}

Longitudinal pitch stability is dictated by the angular inclination of the wishbone pivot axes relative to the chassis centerline. 

To combat severe forward weight transfer during 5G braking zones, the front suspension employs 25-35\% Anti-Dive geometry. The theoretical anti-dive percentage ($AD_{\%}$) is derived mathematically from the side-view geometry:

\begin{equation} \label{eq:anti_dive}
\begin{split}
AD_{\%} &= \left( \frac{\tan(\alpha)}{\tan(\theta_{bias})} \right) \cdot 100
\end{split}
\end{equation}

where $\alpha$ is the angular inclination of the front wishbone reaction vector, and $\theta_{bias}$ represents the angular vector of the mandated braking bias. This geometry mechanically converts the horizontal braking torque into a vertical lifting force, physically resisting chassis nose dive without relying on ultra-stiff, non-compliant springs.

Symmetrically, 40-50\% Anti-Squat geometry is utilized on the rear axle. This is utterly critical for the 2026 regulations; the rear anti-squat geometry physically prevents the rear chassis from bottoming out under the instantaneous, massive torque delivery of the 350 kW MGU-K during aggressive corner exits. However, a delicate tradeoff exists: excessive anti-squat mechanically feeds massive, rigid jacking forces directly into the chassis footprint, permanently disrupting the aerodynamic airflow through the sensitive 30 mm rear diffuser volume.

\subsection{Ride Frequency}

The primary vertical stiffness of a given axle is quantified by its undamped natural ride frequency ($f$):

\begin{equation} \label{eq:ride_frequency}
\begin{split}
f &= \left(\frac{1}{2\pi}\right) \cdot \sqrt{\frac{k}{m_{corner}}}
\end{split}
\end{equation}

where $k$ is the effective wheel-rate spring stiffness and $m_{corner}$ constitutes the sprung mass resting on that specific corner. 

To contextualize the hostility of an F1 platform, a standard commercial road vehicle operates between 1-2 Hz. The 2026 vehicle fundamentally demands a front ride frequency of 6-8 Hz. This hyper-stiff front baseline prioritizes absolute aerodynamic platform consistency, locking the front wing element immovably at its 32 mm ride height target regardless of braking forces. The rear axle utilizes a marginally softer 5-7 Hz target to permit slight mechanical compliance, prioritizing longitudinal traction elasticity for the MGU-K deployment. To decouple ride stiffness from roll stiffness, a third cross-linked element (the heave spring) independently controls purely vertical bilateral motion without interfering in single-wheel bump characteristics.
\end{multicols}

\begin{figure}[H]
\centering
\includegraphics[width=\textwidth]{"figure12_suspension.png"}
\caption{Suspension Geometry Block Diagram mapping Wishbone Layouts, Kinematic Instant Centers, Pullrod Actuation, and Inboard Rocker Assemblies.}
\label{fig:suspension_geometry}
\end{figure}

\begin{multicols}{2}
\subsection{Inerter (J-Damper) Architecture}

Working in strict parallel with the standard torsion springs and rotary hydraulic dampers, the suspension relies on an Inerter (historically known as the J-Damper) \cite{inerter_mechanics_lit}. Unlike a spring (which resists pure displacement) or a damper (which resists pure velocity), an inerter strictly resists physical suspension acceleration. 

The resistive force ($F_{inerter}$) is defined as:

\begin{equation} \label{eq:inerter_force}
F_{inerter} = b \cdot \ddot{x}
\end{equation}

where $b$ is the calibrated inertance (expressed in kg), and $\ddot{x}$ is the linear acceleration of the suspension pushrod. The inerter mechanically behaves as virtual mass added to the corner assembly—smoothing out high-frequency transients without imparting the devastating physical mass penalty of actual ballast, critical for the tight 800 kg limit.

Combining the spring, damper, and inerter yields the total continuous-time Laplace suspension transfer function ($H(s)$):

\begin{equation} \label{eq:transfer_function}
\begin{split}
H(s) &= \frac{bs^2 + c_d s + k_s}{m_s s^2 + bs^2 + c_d s + k_s}
\end{split}
\end{equation}

where $c_d$ is the damping coefficient, $k_s$ is the physical spring rate, and $m_s$ is the sprung corner mass. By precisely tuning the $b$ parameter, vehicle dynamics engineers can aggressively shift the resonant frequency of the tyre-hop mode (typically vibrating violently between 5-20 Hz). Suppressing this specific resonant bandwidth is absolutely vital for ensuring the physical rubber of the tire remains ceaselessly cemented to the tarmac, actively extracting maximum mechanical grip throughout the violent operational envelope of the vehicle.
\end{multicols}

\begin{table}[H]
\centering
\caption{2026 Primary Suspension Geometry and Vehicle Dynamics Parameterization}
\vspace{0.2cm}
\begin{tabularx}{\textwidth}{@{}lll@{}}
\toprule
\textbf{Parameter} & \textbf{Front Axle} & \textbf{Rear Axle} \\
\midrule
Kinematic Configuration & Unequal Double Wishbone & Unequal Double Wishbone \\
Inboard Actuation & Pullrod & Pullrod \\
Track Width & 1600 mm & 1620 mm \\
Static Camber Range & -2.5$^\circ$ to -3.5$^\circ$ & -1.5$^\circ$ to -2.0$^\circ$ \\
Static Toe & -0.05$^\circ$ (divergent out) & +0.1$^\circ$ (convergent in) \\
Caster Angle & 10-12$^\circ$ & --- \\
Kingpin Inclination (KPI) & 8-10$^\circ$ & --- \\
Pitch Resistance Geometry & 25-35\% Anti-Dive & 40-50\% Anti-Squat \\
Undamped Ride Frequency & 6-8 Hz & 5-7 Hz \\
\bottomrule
\end{tabularx}
\label{tab:suspension_geometry}
\end{table}
